\lecture{21}
\title{Quaternions, $O_4(\mathbb{R})$, $SO_4(\mathbb{R})$, $SO_3(\mathbb{R})$, and Conjugacy Classes of $SU_2$}

\begin{document}

\subsection{Quaternion Inner Product and $\{O_4(\mathbb{R}), SO_4(\mathbb{R}), SO_3(\mathbb{R})\}$}

Let's continue our discussion of the quaternions $\mathbb{H} = \mathbb{R} \oplus \mathbb{R}i \oplus \mathbb{R}j \oplus \mathbb{R}k$.

\begin{remark}
    Think of $\mathbb{H}$ as the vector space $\mathbb{R}^4$ with the additional operation ``multiplying two vectors''.
\end{remark}

\begin{theorem}{Commutative Real Parts}
    For any $\alpha, \beta \in \mathbb{H}$, we have $\Real[\alpha\beta] = \Real[\beta\alpha]$.
\end{theorem}

\begin{proof}
    Say $\alpha = \alpha_1 + \alpha_2i + \alpha_3j + \alpha_4k$ and $\beta = \beta_1 + \beta_2i + \beta_3j + \beta_4k$.  Then just compute. ~ $\blacksquare$
\end{proof}

\textbf{Definition (Quaternion Inner Product).} Given $\alpha = \alpha_1 + \alpha_2i + \alpha_3j + \alpha_4k$ and $\beta = \beta_1 + \beta_2i + \beta_3j + \beta_4k$, their \textbf{inner product} is $\langle \alpha, \beta \rangle := \alpha_1\beta_1 + \alpha_2\beta_2 + \alpha_3\beta_3 + \alpha_4\beta_4 = \mathrm{Re}[\alpha \overline{\beta}]$.  (Check these are equivalent!)

So $\mathbb{H}$ is a four-dimensional Euclidean space.  This yields a stronger sense in which $\mathbb{H} \cong \mathbb{R}^4$.  In particular\dots

\begin{boxed-text}[32em]
    \textbf{Observation.} The group of inner-product-preserving linear maps $f: \mathbb{H} \to \mathbb{H}$ \\ is isomorphic to $O_4(\mathbb{R})$ via the usual isomorphism $\mathbb{H} \leftrightarrow \mathbb{R}^4$. \vskip{0.5em}

    In particular, $\{1, i, j, k\}$ is an orthonormal basis of $\mathbb{H}$ under this inner product.
\end{boxed-text}

To classify these ``orthogonal'' linear maps $f: \mathbb{H} \to \mathbb{H}$, we'll need to work with $\mathbb{H}$ and $O_4(\mathbb{R})$ simultaneously!

\begin{theorem}{Quaternions in $O_4(\mathbb{R})$}
    The only inner-product-preserving linear maps $f: \mathbb{H} \to \mathbb{H}$ are $f: x \mapsto \alpha x \beta$ or $f : x \mapsto \alpha \overline{x} \beta$, for $\alpha, \beta \in \mathbb{H}$ satisfying $|\alpha| = |\beta| = 1$.
\end{theorem}

\begin{proof}
    To show the proposed $f: x \mapsto \alpha x \beta$ works, it suffices to show $f$ maps $\{1, i, j, k\}$ to an orthonormal basis.
    \[ \langle f(1), f(i) \rangle = \langle \alpha \beta, \alpha i \beta \rangle = \Real[\alpha\beta  \overline{\alpha i \beta} ] = \Real[\alpha \beta \overline{\beta}(-i) \overline{\alpha}] = \Real [ \overline{\alpha} \alpha (-i)] = 0 ~ \text{ since } \overline{\alpha}\alpha \in \mathbb{R}. ~~~~~ \checkmark \]
    (We use $\Real[\alpha \beta ] = \Real [ \beta \alpha ]$ a lot here!) Checking $f: x \mapsto \alpha \overline{x} \beta$ works analogously.

    To show these $f$ account for \emph{all} inner-product-preserving linear maps, we'll need the $\mathbb{H} \leftrightarrow \mathbb{R}^4$ observation.
    \begin{quote}
        \textbf{Claim.} All of $O_4(\mathbb{R})$ is generated by reflections.

        \emph{Proof:} From Problem Set \#6, we know every element in $O_4(\mathbb{R})$ is the direct sum of two $2 \times 2$ orthogonal submatrices.  And $2 \times 2$ orthogonal matrices are generated by reflections. ~ $\square$
    \end{quote}
    It remains to prove the following about reflections:
    \begin{quote}
        \textbf{Claim.} Every reflection in $\mathbb{R}^4$ is of the form $f_{\alpha}: x \mapsto - \alpha \overline{x} \alpha$ for some $|\alpha| = 1$.

        \emph{Proof:} We claim $f_{\alpha}$ is the reflection sending the basis $\{\alpha, \alpha i, \alpha j, \alpha k\}$ to $\{- \alpha, \alpha i, \alpha j, \alpha k\}$.  Indeed,
        \[ f_{\alpha}(\alpha i) = - \alpha \overline{\alpha i} \alpha = - \alpha (-i) \overline{\alpha} \alpha = \alpha i |\alpha |^2 = \alpha i ~~~ \text{ and } ~~~ f_{\alpha}(\alpha) = - \alpha \overline{\alpha} \alpha = - |\alpha|^2 \alpha = -\alpha. \]
        So $f_{\alpha}$ is a reflection about the hyperplane perpendicular to $\alpha$. ~ $\square$
    \end{quote}
    Thus, our proposed maps $f: x \mapsto \alpha x \beta$ and $f: x \mapsto \alpha \overline{x} \beta$ include all reflections, which generate $O_4(\mathbb{R})$; and they are closed under composition, so they are \emph{all} of $O_4(\mathbb{R})$. ~ $\blacksquare$
\end{proof}

\begin{remark}
    Multiplying by elements of $\mathbb{H}$ yields ``rotations'' of $\mathbb{R}^4$, whereas conjugation yields ``reflections'' of $\mathbb{R}^4$.
\end{remark}

In summary, recalling that $\{\alpha \in \mathbb{H} : |\alpha| = 1 \} \cong S^3$, we have the following description of $O_4(\mathbb{R})$.
\[ \varphi: S^3 \times S^3 \to O_4(\mathbb{R}) ~ \text{ via } ~ \varphi: (\alpha, \beta) \mapsto \begin{cases}
    \,  (x \mapsto \alpha x \beta) \\
    \, (x \mapsto \alpha\overline{x} \beta)
\end{cases} \]
More specifically, one can restrict the above to describe only maps in $SO_4(\mathbb{R})$.
\[ \varphi: S^3 \times S^3 \to SO_4(\mathbb{R}) ~ \text{ via } ~ \varphi: (\alpha, \beta) \mapsto (x \mapsto \alpha x \beta) ~ \text{ with kernel } ~ \ker \varphi = \{ \pm (1, 1)\}. \]
This makes sense---conjugations yield ``reflections'', which have negative determinant.  By corollary\dots

\begin{theorem}{Describing $SO_3(\mathbb{R})$}
    All rotations of $\mathbb{R} i \oplus \mathbb{R}j \oplus \mathbb{R}k$ are $f: x \mapsto \alpha x \alpha^{-1}$ for $\alpha \in \mathbb{H}$ with $|\alpha| = 1$.
\end{theorem}

\begin{proof}
    Think of $SO_3(\mathbb{R})$ as the subset of $SO_4(\mathbb{R})$ that fixes $1$.  The only maps $(x \mapsto \alpha x \beta)$ that fix $1$ are those for which $\alpha \beta = 1$, or $\beta = \alpha^{-1}$.  This yields:
    \[ \varphi:  S^3 \to SO_3(\mathbb{R}) ~ \text{ via } ~ \varphi : \alpha \mapsto (x \mapsto \alpha x \alpha^{-1}) ~ \text{ with kernel } ~ \ker (\varphi) = \{ \pm 1\}. \]
    And that's what we wanted. ~ $\blacksquare$
\end{proof}

This is more of a ``real-life'' application---it implies the following:

\begin{boxed-text}[33em]
\textbf{Claim.} Three-dimensional rotations are parametrized by $S^3 / \{ \pm 1\}$. \vskip{0.5em}

\emph{Parametrization:} For any $\alpha \in S^3$, we have a corresponding $\alpha \in \mathbb{H}$ with $|\alpha| = 1$. \\ This yields the $\mathbb{H} \to \mathbb{H}$ map defined via $x \mapsto \alpha x \alpha^{-1}$. Then taking the $(i, j, k)$ \\ components of this $\mathbb{H} \to \mathbb{H}$ map yields a well-defined $\mathbb{R}^3 \to \mathbb{R}^3$ map in $SO_3(\mathbb{R})$.
\end{boxed-text}

Three-dimensional computer graphics use quaternions for this reason!

\subsection{Conjugacy Classes of $SU_2$}

Recall from last lecture the following interpretations of $SU_2$.
\begin{align*}
SU_2 \cong \ & \left \{ \begin{bsmallmatrix} \alpha & \beta \\ - \overline{\beta} & \overline{\alpha} \end{bsmallmatrix} : \alpha, \beta \in \mathbb{C} \text{ and } |\alpha|^2 + |\beta|^2 = 1\right \} \cong  \{v \in \mathbb{R}^4: |v| = 1 \} \cong \{ \alpha \in \mathbb{H} : |\alpha| = 1 \}\cong S^3. 
\end{align*}
To be more precise about these isomorphisms\dots
\begin{itemize}
    \item The isomorphism $SU_2 \cong \{v \in \mathbb{R}^4 : |v| = 1\}$ is via $\begin{bsmallmatrix} \alpha & \beta \\ -\overline{\beta} & \overline{\alpha} \end{bsmallmatrix} \mapsto (\mathrm{Re}[\alpha], \mathrm{Im}[\alpha], \mathrm{Re}[\beta], \mathrm{Im}[\beta]) \in \mathbb{R}^4$.
    \item The isomorphism $SU_2 \cong \{ z \in \mathbb{H} : |z| = 1\}$ is via $\begin{bsmallmatrix} \alpha & \beta \\ -\overline{\beta} & \overline{\alpha} \end{bsmallmatrix} \mapsto \alpha + \beta j \in \mathbb{H}$.
\end{itemize}
Since $SU_2 \cong S^3$, the previous section showed $SU_2 / \{\pm 1\} \cong SO_3(\mathbb{R})$; that is, $SU_2$ is a \emph{double cover} of $SO_3(\mathbb{R})$.

\begin{theorem}{Conjugacy Classes of $SU_2$}
    Two matrices in $SU_2$ are conjugate iff they have the same trace.
\end{theorem}

\begin{proof}
    To show conjugates always have the same trace, just recall the identity $\mathrm{tr}(XY) = \mathrm{tr}(YX)$ and compute:
    \[ \mathrm{tr}(ABA^{-1}) = \mathrm{tr}(A^{-1}AB) = \mathrm{tr}(B). \]
    (More strongly: conjugate matrices always have the same spectra because they're a change-of-basis apart!)

    For the other direction, pick some constant $c$, and consider \emph{any} $A \in SU_2$ with $\mathrm{tr}(A) = c$.  We wish to show that the condition $\mathrm{tr}(A) = c$ uniquely determines $A$ up to conjugation in $SU_2$.
    
    By the Spectral Theorem, we have $A = M \Lambda M^{-1}$ for some diagonal matrix $\Lambda$ and some $M \in U_2$.  Well\dots
    \begin{itemize}
        \item Taking $M' := \frac{M}{(\det M)^{1/2}}$, we see $A = M' \Lambda (M')^{-1}$ for some $M' \in SU_2$.
        \item Furthermore, we have the following constraints on $\Lambda = \begin{bsmallmatrix} \lambda_1 & 0 \\ 0 & \lambda_2 \end{bsmallmatrix}$.
    \[ \det A = 1 \implies \lambda_1 \lambda_2 = 1 ~~~~~ \text{ and } ~~~~~ \mathrm{tr}(A) = c \implies \lambda_1 + \lambda_2 = c \]
    Thus, $\{\lambda_1, \lambda_2\}$ is the set of roots $\{r_1, r_2\}$ of $X^2 - cX + 1 = 0$.  In other words, either $\Lambda = \begin{bsmallmatrix} r_1 & 0 \\ 0 & r_2 \end{bsmallmatrix}$ or $\Lambda = \begin{bsmallmatrix} r_2 & 0 \\ 0 & r_1 \end{bsmallmatrix}$.
    \end{itemize}
    Therefore, \emph{any} $A \in SU_2$ with $\mathrm{tr}(A) = c$ is conjugate in $SU_2$ to either $\begin{bsmallmatrix} r_1 & 0 \\ 0 & r_2 \end{bsmallmatrix}$ or $\begin{bsmallmatrix} r_2 & 0 \\ 0 & r_1 \end{bsmallmatrix}$ for fixed constants $\{r_1, r_2\}$.  And these two matrices are conjugates, too (obviously!):
    \[ \begin{bmatrix} 0 & -1 \\ 1 & 0 \end{bmatrix} \begin{bmatrix} r_1 & 0 \\ 0 & r_2 \end{bmatrix} \begin{bmatrix} 0 & -1 \\ 1 & 0 \end{bmatrix}^{-1} = \begin{bmatrix} 0 & -1 \\ 1 & 0 \end{bmatrix} \begin{bmatrix} r_1 & 0 \\ 0 & r_2 \end{bmatrix} \begin{bmatrix} 0 & 1 \\ -1 & 0 \end{bmatrix} = \begin{bmatrix} 0 & -1 \\ 1 & 0 \end{bmatrix} \begin{bmatrix} 0 & r_1 \\ -r_2 & 0 \end{bmatrix} = \begin{bmatrix} r_2 & 0 \\ 0 & r_1 \end{bmatrix}. \]
    And so we win. ~ $\blacksquare$
\end{proof}

\begin{remark}
    (Reading Comprehension) In the last step, why did we conjugate by $\begin{bsmallmatrix} 0 & -1 \\ 1 & 0 \end{bsmallmatrix}$ instead of $\begin{bsmallmatrix} 0 & 1 \\ 1 & 0 \end{bsmallmatrix}$?
\end{remark}

\subsection{Geometry of $SU_2$: Latitude and Longitude}
    
Trace in $SU_2$ has a nice analogue in $\mathbb{H}$, which you can check by plain computation:

\begin{theorem}{Redefining Trace}
    Every $A = \begin{bsmallmatrix} \alpha & \beta \\ -\overline{\beta} & \overline{\alpha} \end{bsmallmatrix} \in SU_2$ yields a $\gamma = \alpha + \beta j \in \mathbb{H}$.  Then $\mathrm{tr}(A) = 2 \mathrm{Re}[\gamma]$.
\end{theorem}

This inspires the following \emph{geometric} description of conjugacy classes in $SU_2$.
\begin{center}
    \frame{\includegraphics[width=0.25\textwidth]{lecture-21/latitude-classes}}
\end{center}
The diagram above has the following features:
\begin{itemize}
    \item The space $SU_2 \cong \{ \alpha \in \mathbb{H} : |\alpha| = 1\} \cong S^3$ is depicted as the blue surface of a sphere.
    \begin{itemize}
        \item Technically speaking, this is wrong!  Notably, $S^3$ is the surface of the four-dimensional ball $B^4$.  But we can't draw $B^4$, so we'll have to settle for using $S^2$ to visually represent $S^3$.
    \end{itemize}
    \item Conjugacy classes---that is, collections of matrices with the same trace---look like green lines of latitude.
\end{itemize}

\begin{remark}
Keep in mind that everything we say about $SU_2$ applies equally well to the subgroup $\{\alpha \in \mathbb{H} : |\alpha | = 1 \}$ of $\mathbb{H}$.  In particular, multiplication in $SU_2$ is isomorphic to multiplication in $\mathbb{H}$.
\begin{itemize}
    \item (Reading Comprehension) Convince yourself that the North Pole is $\{+I_2\}$, and the South Pole is $\{-I_2\}$.
    \item (Reading Comprehension) Convince yourself that the remaining conjugacy classes of $SU_2$ are isomorphic to $S^2$.
    \item (Reading Comprehension) Convince yourself that $i$, $j$, and $k$ live on the equator of $S^3$.
\end{itemize}
\end{remark}


We've interpreted the lines of latitude---what about the lines of longitude?
\begin{center}
    \frame{\includegraphics[width=0.4\textwidth]{lecture-21/longitude-subgroups}}
\end{center}
The diagram above has the following features:
\begin{itemize}
    \item The red equator meets the orange line of longitude at some point $x \in S^3$ with $\mathrm{Re}[x] = 0$. \\ (Or equivalently, $A \in SU_2$ with $\mathrm{tr}(A) = 0$.)
    \item The line of longitude is parametrized by $\{ 1 \cdot \cos \theta + x \cdot \sin \theta \mid \theta \in [0, 2\pi) \}$.
\end{itemize}

\begin{theorem}{Longitudinal Subgroups}
    For any $x \in \{ \alpha \in \mathbb{H} : |\alpha| = 1\}$ with $\mathrm{Re}[x] = 0$, the line of longitude $H_x := \{1 \cdot \cos \theta + x \cdot \sin \theta \mid \theta \in [0, 2\pi)\}$ is a subgroup of $\mathbb{H}$.
\end{theorem}

\begin{proof}
    Both $x \in \mathbb{H}$ and $i \in \mathbb{H}$ lie on the equator, but the equator is a conjugacy class!  So $x$ and $i$ are conjugates:
    \[ H_x := \{1 \cdot \cos \theta + x \cdot \sin \theta \mid \theta \in [0, 2\pi)\} ~ \text{ is conjugate to } ~ H_i := \{1 \cdot \cos \theta + i \cdot \sin \theta \mid \theta \in [0, 2\pi)\} \]
    But $H_i$ is obviously a subgroup of $\mathbb{H}$---it's literally just $\{z \in \mathbb{C} : |z| = 1\}$. ~ $\blacksquare$
\end{proof}

\subsection{Geometry of $SU_2$: Normal Subgroups and Simplicity}

One last thing: we'll use the results we've established so far to provide a geometrical proof of the following fact!

\begin{boxed-text}[29em]
    \textbf{Claim.} The only normal subgroups of $SU_2$ are $\{1\}$, $\{1, -1\}$, and $SU_2$.
\end{boxed-text}

Say $G$ is a normal subgroup of $SU_2$ that is not either $\{1\}$ or $\{1, -1\}$.  We wish to show $G = SU_2$.

We may pick some $x \in G$ such that $x \not \in \{1, -1\}$, marked in the diagram below.

\begin{center}
    \frame{\includegraphics[width=0.2\textwidth]{lecture-21/sphere-normality}}
\end{center}

The driving principle behind our proof will be the following fact:
\begin{quote}
    \textbf{Fact.} Since $G$ is normal, it is the union of conjugacy classes of $SU_2$.
\end{quote}
Let $S$ be the conjugacy class of $x$, shown as a green latitudinal line.
\begin{itemize}
    \item Since $x \in G$, and $G$ is normal, we must also have $S \subseteq G$.
    \item By closure of groups, we also have $x^{-1}S \subseteq G$, shown in orange.
    \item Every conjugacy class (highlighted in green) that intersects with an element of $x^{-1}S$ must also be in $G$.
\end{itemize}
Thus, just from assuming $x \in G$, we've shown $G$ must also contain the entire region highlighted in green!

Now consider any longitudinal subgroup $H$ of $SU_2$.  Since $G$ contains the whole region highlighted in green, $G$ must also contain an entire arc of $H$.  And an arc generates all of $H$, so $H \subseteq G$.

But the longitudinal subgroups of $SU_2$ together cover all of $SU_2$.  So $G$ contains all of $SU_2$, done. ~ $\blacksquare$

\begin{remark}
    By corollary, $SO_3(\mathbb{R}) \cong SU_2 / \{ \pm 1 \}$ is simple.
\end{remark}

\end{document}